% Punto 2 · Metodología.
% Texto de la compañera, tomado de codigo/punto2/00_generar_html.py (apartados 2.1 a 2.5, párrafos de método) y
% pasado a LaTeX tal cual. Cambios: formato, punto decimal y espacio para los miles (18.000 → 18 000; 3,5 → 3.5),
% los títulos de apartado pasan a \subsubsection sin numeración del taller, y la justificación del tc es la que
% Juan decidió el 26 nov (columna tc_selection_note de productos/punto2/problemaA_alcantarilla.csv).
% Las cifras se revisaron contra los CSV de productos/punto2/ solo en parte (26 nov); falta pasarlas a macros.

\subsubsection{Datos subdiarios y control de calidad}
Cada año se acepta si cumple simultáneamente $\geq 80$~\% de horas válidas en el año y $\geq 90$~\% entre abril y
septiembre; el resto se descarta. La segunda regla se detalla en el archivo anual
(\path{productos/punto2/resumen_completitud_anual.csv}).

Diez eventos horarios mayores: revisión de acumulados. Las ventanas previas y posteriores reportan suma de lluvia y
horas válidas dentro de 24~h; los códigos conservan el contexto del registro. Esta inspección detecta picos cercanos
a vacíos, pero la confirmación meteorológica requiere contrastar reportes y fuente original.

Sensibilidad al filtro de completitud: comparación de máximos anuales horarios y cuantiles preliminares $T=10$ y
$T=50$. Las estimaciones se calcularon con Gumbel por momentos; la muestra de años y el supuesto de distribución
limitan su interpretación.

\paragraph{Resolución y uso en microcuencas.} La serie tiene resolución horaria. No permite medir directamente máximos
de 10, 15 o 30 minutos, que pueden controlar la respuesta de microcuencas. La corrección de \citet{weiss1964} ajusta
máximos obtenidos con reloj fijo frente a máximos de duración móvil y se aplica a la duración observada; no
reconstruye variabilidad subhoraria. Para duraciones menores de una hora se necesita una ley de escala
intensidad--duración (idealmente regional y validada), además de cuantificar su incertidumbre.

\subsubsection{Máximos, distribuciones e incertidumbre}
Se estiman máximos móviles por año hidrológico octubre--septiembre. Las intensidades incluyen la corrección de Weiss
(1964), con $f_W = 1/(1 - 1/(8n))$; la columna de coherencia verifica que las láminas corregidas no decrezcan al
aumentar la duración. El año hidrológico se identifica por el año en que termina.

MOM: momentos; LMOM: L-momentos muestrales no sesgados según las fórmulas PWM de \citet{hosking1997}; MLE: máxima
verosimilitud. En GEV se informa explícitamente $\xi$ y el parámetro de SciPy $c = -\xi$. Las posiciones de
graficación son las de Gringorten ($a = 0.44$; \citealp{gringorten1963}).

KS y AD comparan la distribución ajustada con los máximos observados; PPCC es la correlación producto-momento en el
gráfico de probabilidad. Mann--Kendall \citep{mann1945} evalúa tendencia monotónica y la pendiente de Sen \citep{sen1968} se expresa en mm/h por año
hidrológico.

\subsubsection{Curvas intensidad--duración--frecuencia}
Se toman los cuantiles de Log-Pearson III ajustado por L-momentos (LP3/LMOM) como base y se comparan con GEV y Gumbel.
Se impuso coherencia física con regresión isotónica sobre $P(d,T)=i\cdot d$ para hacer no decreciente la lámina con
duración y sobre la intensidad frente al periodo de retorno para evitar cruces. Las curvas paramétricas tienen
exponente temporal $m$ positivo. Las bandas subhorarias se obtienen ajustando los límites bootstrap por duración con
el modelo IDF elegido.

\paragraph{Selección.} Bernard/Sherman, Talbot/Wenzel y Koutsoyiannis \citep{koutsoyiannis1998} se compararon con MAPE, RMSE y $R^2$. La tabla
marca el mejor ajuste global y, por separado, el mejor de los modelos con desplazamiento $b>0$ utilizado para la tabla
maestra y la extrapolación a duraciones cortas, lo que limita la intensidad cuando $d$ tiende a cero.

\paragraph{Validación con NOAA Atlas 14.} Se cotejan las 36 combinaciones de duración y periodo de retorno con las
profundidades PDS puntuales obtenidas del endpoint oficial NOAA HDSC para Denver Stapleton \citep{perica2013}. La diferencia relativa
compara nuestra IDF LP3/LMOM--Talbot/Wenzel con NOAA.

\paragraph{Periodo de retorno de eventos históricos.} Los periodos se estiman con LP3/LMOM sobre la serie de máximos
anuales aceptada. Para el evento de 2022 se usa un umbral espacial/subhorario documentado \citep{accuweather2022}, por lo que el resultado es
una cota inferior y no un periodo puntual observado en la estación.

\subsubsection{Diseño hidrológico de la microcuenca}
La sección separa la cuenca topográfica delimitada con DEM del escenario metodológico NRCS que se conserva como
respaldo. El DEM Copernicus es un modelo de superficie y la cuenca automática requiere validación de drenaje urbano
antes de usarla para construcción. La cuenca candidata y el cauce se extraen del DEM Copernicus GLO-30 \citep{esa_copdem} reproyectado a
UTM 13N.

\paragraph{Problema A: alcantarilla circular.} Predimensionamiento con geometría leída de los metadatos DEM, cobertura
periurbana ponderada, IDF LP3/LMOM y Talbot--Wenzel, Manning a tubo lleno y diámetro comercial superior. El diseño usa
$T=50$ años y contrasta el criterio de riesgo con \citet{chow1988}. El $t_c$ de diseño es 60~min: Kerby--Hathaway
(69.8; \citealp{txdot_hdm}), Témez (82.4; en la forma de la norma 5.2-IC, \citealp{fom2016}) y NRCS-SCS Lag (175.0~min; \citealp{nrcs2010cap15}) quedan por encima; Kirpich (45.5~min; \citealp{txdot_hdm}) es el valor bajo porque solo
representa el tránsito por el cauce y omite el flujo laminar inicial.

\paragraph{Problema B: tanque de detención por urbanización.} El escenario cambia de $C=0.495$ y $t_c=60$~min a
$C=0.78$ y $t_c=30$~min, manteniendo $T=50$ años y limitando la descarga a $Q_\mathrm{lim}=Q_{p,\mathrm{previo}}=
6.426$~m$^3$/s. El tanque prismático tiene huella 18\,000~m$^2$, altura útil 3.5~m y volumen geométrico
63\,000~m$^3$; la salida se controla con orificio dimensionado ($C_d = 0.61$) y tránsito RK4 de 1~min
($S = A_\mathrm{tanque}\,h$).

\paragraph{Respaldo metodológico: hietogramas NRCS-UH y AMC.} Hietogramas de bloques alternos e hidrogramas NRCS-UH \citep{nrcs2007cap16}
de la misma microcuenca DEM y con los mismos parámetros del Problema A ($A=1.0269$~km$^2$, $C=0.495$,
$CN_\mathrm{II}=78$ y $t_c=60$~min). El caudal racional de esta sección usa el mismo $C$, sin factor de frecuencia
($C_f = 1$), así que coincide con el del Problema A.

\subsubsection{Análisis complementarios de diseño y riesgo}
Cuatro análisis complementarios cuantifican el compromiso entre costo de protección, extremos por excedencias,
calentamiento futuro y desempeño histórico de la alcantarilla. Los costos se expresan en unidades de $C_0$
normalizadas porque los archivos fuente no especifican un costo monetario base.

\paragraph{Costo esperado de protección.} Se evalúan $T_d=2$--500 años con $Q_d=Q(1/T_d)$, costo anualizado
$\mathrm{FRC}\cdot C_\mathrm{obra}$ y daño anual esperado integrado entre $p=0$ y $1/T_d$.

\paragraph{POT/GPD y relación de Langbein.} La serie horaria se declusterizó con separación mínima de seis horas
secas. El umbral seleccionado genera alrededor de dos eventos por año válido; se ajustó GPD \citep{coles2001} a los excesos con
convención scipy $c=+\xi$ y se transformó el retorno de evento a retorno AMS mediante \citet{langbein1949}.

\paragraph{Cambio climático y GEV no estacionaria.} La curva IDF se escala con 7~\% por °C, la tasa de Clausius--Clapeyron \citep{fowler2021}. La GEV evalúa
$\mu(t)=\mu_0+\mu_1 t$ en $t_\mathrm{final}+25$ años contra el umbral fijo $i^*=48.93$~mm/h.

\paragraph{Confiabilidad, resiliencia y vulnerabilidad.} Hashimoto \citep{hashimoto1982} se calculó con caudal horario, Weiss $=1.14$ y
capacidades a tubo lleno de Manning. $\alpha$ es la fracción de horas sin falla; $\gamma$ es la tasa de recuperación
por hora fallida; $\nu$ resume el exceso sobre capacidad.
